\subsubsection{Subproceso 3.B: Confirmación y Registro Transaccional Oficial}

\begin{enumerate}
    \item \textbf{Asentamiento transaccional definitivo de la carga aprobada.} \\
    \textit{Responsable: Intranet DIRCE / Sistema SIGA.} \\
    Al recibir la Pre-Ficha Aprobada o Subsanada del Subproceso 3.A, el sistema bloquea la fase de edición del estudiante y graba las cargas lectivas en la base de datos central, dotando al registro de carácter legal y transaccional.

    \item \textbf{Generación del Reporte de Matrícula.} \\
    \textit{Responsable: Sistema SIGA.} \\
    El sistema emite el documento digital preliminar que acredita las asignaturas registradas para el ciclo en curso, en concordancia con la condición de aprobación automática del procedimiento (TUPA PA5400954D, RR-3698-2024).

    \item \textbf{Cierre de Matrícula por Facultad.} \\
    \textit{Responsable: Oficina de Estadística de la Facultad.} \\
    Tras la semana de depuración —en la que recaba la estadística por sección (RR-0570-2022, Art. 75°) y tramita las rectificaciones autorizadas por la Dirección de Escuela (RR-0570-2022, Art. 77°)—, la oficina consolida el padrón, indica a la DIRCE la culminación del proceso y remite el archivo digital oficial de consolidación (RR-0570-2022, Art. 78°).

    \item \textbf{Emisión y visado de la Boleta de Matrícula Oficial.} \\
    \textit{Responsable: DIRCE / Oficina de Estadística de la Facultad.} \\
    La DIRCE genera el registro legal de la Boleta en el SIGA y la Oficina de Estadística aplica el visado institucional que otorga plena validez jurídica al Reporte y a la Boleta.

    \item \textbf{Generación y distribución de padrones de clase.} \\
    \textit{Responsable: Oficina de Estadística de la Facultad / DIRCE.} \\
    Se preparan los listados oficiales definitivos por asignatura y sección, y se ponen a disposición de cada docente el primer día hábil de la segunda semana lectiva (RR-0570-2022, Art. 78°).

    \item \textbf{Reporte institucional y tramitación externa.} \\
    \textit{Responsable: DIRCE.} \\
    La dependencia migra los datos oficiales a los sistemas de SUNEDU y MINEDU en el plazo de veinte días (RR-0570-2022, Art. 79°), habilita el trámite de expedición del carné universitario y, en un plazo no mayor a diez días hábiles tras el cierre, pone a disposición de las Facultades el registro de matrícula, las boletas y el SIGA actualizado (RR-0570-2022, Art. 80°).
\end{enumerate}
